\chapter{解题感悟}
\label{cp:notes}
% TODO: 在这里写第三章内容，替换掉示例文字。

%In the course of writing, it is sometimes necessary to insert notes in the middle of the document so that in the future we can return to them and read what we have written. Sometimes we need a space to jot down a quick \verb|TODO|. With this in mind, two commands have been developed: \verb|\todo{TEXT}| and \verb|\note{TEXT}|. These commands accept normal text and produce two different types of boxes, respectively, for each type of use. These boxes are ``breakable'', which means that they can be placed anywhere in the document and will behave as if they were text. Below is a simple usage of both environments. Please note that both of these environments were created with the intention of helping you develop your work and \textbf{should not be used in the final version of your document!}

%\todo{Write about X! Later, we should remove Y.}

%\note{Come to me, all you who are weary and burdened, and I will give you rest. \citep{一个人类}. Take my yoke upon you and learn from me, for I am gentle and humble in heart, and you will find rest for your souls.}


这周在经过上周的练习后，对命令行操作这些流程更加熟悉了。解决问题也没有上次那么痛苦了。虽然还是有一些问题需要查阅资料，但至少比上一次清楚了很多。

然后因为上次已经配置好了latex环境，学习了怎么使用latex写作，这次实验报告的撰写也明显顺利了许多，完全没有上次那么痛苦了。

希望经过这几次的练习，下次完成练习能更加顺利，以后也能更加熟练的使用命令行操作和latex写作。